\chapter{Függelékek és production checklist}

Ez a fejezet nem új nyelvi elemeket vezet be. Az előző fejezetekből összeálló, röviden használható üzemeltetési segédlet: tipikus könyvtárstruktúra, reverse proxy minták, systemd unit, logrotate és egy release előtti ellenőrzőlista. A minták kiindulópontok; a konkrét domaineket, tanúsítványokat, jogosultságokat és erőforráslimiteket a saját környezethez kell igazítani.

\section{Ajánlott production könyvtárstruktúra}

Az RWLang helyben telepített szoftver, ezért a könyv következetesen a \code{/usr/local} hierarchiát használja a programhoz és a konfigurációhoz. Az alkalmazás állapota és a logok ettől külön élnek.

\begin{lstlisting}
/usr/local/bin/
  rwlang-server
  rwlang-cli

/usr/local/etc/rwlang/
  server.toml
  rate-limits.toml
  resource-profiles.toml

/run/secrets/rwlang/
  database-url
  redis-url
  local-auth-database-url

/srv/rwlang/
  current/                 # current application release
  releases/
    2026-09-04.1/
  data/                    # AppFs / durable application data

/var/log/rwlang/
  server.log
  access.log
  audit.log
\end{lstlisting}

A release tartalom legyen lehetőleg read-only. A szervernek csak azokon a helyeken legyen írási joga, ahol valóban tartós állapotot vagy logot kell írnia. A secret fájl ne legyen az alkalmazás release könyvtárában és ne kerüljön verziókezelésbe.

\section{Apache reverse proxy teljes minimálminta}

A példa egyetlen Apache reverse proxyt feltételez. A publikus TLS az Apache feladata, az RWLang pedig csak loopbacken figyel.

Az ehhez tartozó RWLang minimum:

\begin{lstlisting}
[server]
listen = "127.0.0.1:8080"
insecure_dev_cookies = false

[tls]
public_host = "www.example.com"

[web]
trusted_proxy_cidrs = ["127.0.0.1/32", "::1/128"]
\end{lstlisting}

Reverse proxy production módban a certificate/key a proxyn marad, de a trusted \code{public_host} kell a Host/Origin pinninghez. A proxy feladata a hiteles \code{X-Forwarded-Proto: https} metadata előállítása.

\begin{lstlisting}
<VirtualHost *:80>
    ServerName www.example.com
    Redirect permanent / https://www.example.com/
</VirtualHost>

<VirtualHost *:443>
    ServerName www.example.com

    SSLEngine on
    SSLCertificateFile \
      /etc/letsencrypt/live/www.example.com/fullchain.pem
    SSLCertificateKeyFile \
      /etc/letsencrypt/live/www.example.com/privkey.pem

    ProxyPreserveHost On
    ProxyPass        / http://127.0.0.1:8080/
    ProxyPassReverse / http://127.0.0.1:8080/

    RequestHeader set X-Forwarded-Proto "https"
</VirtualHost>
\end{lstlisting}

Tipikus Apache modulok: \code{proxy}, \code{proxy_http}, \code{headers} és \code{ssl}. Az RWLang trusted-proxy listája csak a tényleges proxyt foglalja magában. Loopback proxy esetén ez tipikusan \code{127.0.0.1/32} és szükség esetén \code{::1/128}. Internetméretű trusted proxy tartomány hibás konfiguráció.

\section{Nginx reverse proxy teljes minimálminta}

\begin{lstlisting}
server {
    listen 80;
    server_name www.example.com;
    return 301 https://$host$request_uri;
}

server {
    listen 443 ssl;
    server_name www.example.com;

    ssl_certificate
      /etc/letsencrypt/live/www.example.com/fullchain.pem;
    ssl_certificate_key
      /etc/letsencrypt/live/www.example.com/privkey.pem;

    location / {
        proxy_pass http://127.0.0.1:8080;
        proxy_http_version 1.1;
        proxy_set_header Host $host;
        proxy_set_header X-Forwarded-Proto https;
        proxy_set_header X-Forwarded-For $remote_addr;
    }
}
\end{lstlisting}

A példában az Nginx közvetlenül látja a klienst, ezért a forwarding címet maga állítja elő. Ne örökíts tovább internet felől érkező tetszőleges \code{X-Forwarded-*} értéket authorityként. Több proxy esetén a trust láncot explicit kell megtervezni.

\section{systemd service minta}

A repository teljes mintája: \code{examples/systemd/rwlang.service}. A lényegi forma:

\begin{lstlisting}
[Unit]
Description=RWLang application
Wants=network-online.target
After=network-online.target

[Service]
Type=simple
User=rwlang
Group=rwlang
WorkingDirectory=/srv/rwlang/current
ExecStartPre=/usr/local/bin/rwlang-server \
  --config /usr/local/etc/rwlang/server.toml \
  --check-config
ExecStart=/usr/local/bin/rwlang-server \
  --config /usr/local/etc/rwlang/server.toml
ExecReload=/bin/kill -HUP $MAINPID
Restart=on-failure
RestartSec=2s
TimeoutStopSec=45s

MemoryMax=1G
MemorySwapMax=0
CPUQuota=200%
TasksMax=256
LimitNOFILE=8192
# Direct TLS 80/443 eseten add hozza:
# AmbientCapabilities=CAP_NET_BIND_SERVICE
# CapabilityBoundingSet=CAP_NET_BIND_SERVICE
NoNewPrivileges=yes
PrivateTmp=yes
PrivateDevices=yes
ProtectSystem=strict
ProtectHome=yes
ProtectKernelTunables=yes
ProtectKernelModules=yes
ProtectControlGroups=yes
RestrictSUIDSGID=yes
LockPersonality=yes
ReadWritePaths=/srv/rwlang/data /var/log/rwlang

[Install]
WantedBy=multi-user.target
\end{lstlisting}

A systemd hard limit és az RWLang request/runtime resource policyja egymást kiegészítő védelmi réteg. A systemd baseline maga legyen a process cgroup authority; ilyenkor ne kapcsold be mellette az RWLang \code{[cgroup]} író módját. Közvetlen 80/443 bindnál csak a \code{CAP_NET_BIND_SERVICE} capabilityt add, reverse proxy mögötti magas portnál ezt is hagyd el. A \code{TimeoutStopSec} legyen nagyobb, mint az alkalmazott graceful shutdown időablak. Behind-proxy módban a SIGHUP tranzakciósan újratöltheti a domain/application hosting állapotot, míg process-szintű konfigurációváltozás után továbbra is kontrollált restart kell; pusztán alkalmazásforrás-változást normál esetben az automatikus source-reload supervisor aktivál.

\section{logrotate minta}

\begin{lstlisting}
/var/log/rwlang/server.log /var/log/rwlang/access.log {
    daily
    rotate 14
    compress
    delaycompress
    missingok
    notifempty
    create 0640 rwlang rwlang
    sharedscripts
    postrotate
        systemctl kill -s HUP --kill-who=main \
          rwlang.service >/dev/null 2>&1 || true
    endscript
}

/var/log/rwlang/audit.log {
    daily
    rotate 90
    compress
    delaycompress
    missingok
    notifempty
    create 0640 rwlang rwlang
    sharedscripts
    postrotate
        systemctl kill -s HUP --kill-who=main \
          rwlang.service >/dev/null 2>&1 || true
    endscript
}
\end{lstlisting}

A konkrét retention jogi és kockázati döntés. Az auditlog általában hosszabb megőrzést és központi továbbítást érdemel; a minta számai nem univerzális előírások.

\section{Release előtti production gate}

A következő lista szándékosan rövid és végrehajtható. A részletes indoklást az előző fejezetek adják.

\begin{enumerate}
  \item \textbf{Forrás és build:} a release commit azonosított; \code{Cargo.lock} rögzített; a platform buildje \code{cargo build --locked --release -p rwlang-server -p rwlang-cli}; az alkalmazás forrása \code{rwlang-cli check} ellenőrzésen átment.
  \item \textbf{Migration:} production előtt \code{migrate status} és \code{migrate verify}; a jóváhagyott schema változás explicit \code{migrate apply}; utána újabb verify.
  \item \textbf{Konfiguráció:} a valódi production \code{server.toml} config-checken átment. Az effective config redaktált kimenete reviewzható.
  \item \textbf{Secrets:} DB/Redis/auth secret fájlok nem a repositoryban vannak, szűk filesystem jogosultsággal rendelkeznek, és nem kerülnek logba vagy release-recordba.
  \item \textbf{HTTP edge:} a domain, TLS és Host kezelés helyes. A trusted-proxy lista kizárólag a tényleges proxykat tartalmazza.
  \item \textbf{Auth:} production cookie policy aktív; local auth/LDAP konfiguráció ellenőrzött; admin/MFA recovery folyamat dokumentált.
  \item \textbf{CSRF és CORS:} state-changing böngészős kérés CSRF-védett; CORS allowlist explicit; credentialed CORS esetén nincs wildcard origin.
  \item \textbf{Adatbázis:} az alkalmazás DB user least privilege; pool/budget összhangban van a backend kapacitásával; fontos queryk indexelése ismert.
  \item \textbf{Storage/upload:} AppFs root confinement és jogosultságok helyesek; upload byte/pixel limitek aktívak; release tartalom nem írható alkalmazásból.
  \item \textbf{Redis/cache/rate limit:} multi-instance környezetben a közös security/runtime state Redisben van; Redis kiesési viselkedése ismert és tesztelt.
  \item \textbf{Egress:} csak named outbound target engedélyezett; host/CIDR/port/TLS/timeout/méret policy explicit; secret nem kerül URL-be vagy logba.
  \item \textbf{Erőforrások:} request/runtime limitek és process hard limitek konzisztensen be vannak állítva; cgroup authorityként systemd vagy külön delegált RWLang cgroup legyen, ne mindkettő.
  \item \textbf{Observability:} server/access/audit log írható; request ID működik; metrics és readiness elérhető a megfelelő belső hálózatról; riasztási minimum definiált.
  \item \textbf{Backup/recovery:} van friss, ellenőrzött DB + AppFs + identity backup; recovery manifest rendelkezésre áll; restore drill nem csak papíron létezik.
  \item \textbf{Deploy/rollback:} a release artifact és checksum ismert; a régi release visszaállítható; schema-változásnál tisztázott, hogy app rollback vagy teljes restore lehetséges.
  \item \textbf{Indítás utáni evidence:} liveness, readiness és legalább egy üzleti smoke kérés sikeres; logban nincs startup fallback vagy váratlan policy-hiba.
\end{enumerate}

\section{Gyors hibadiagnosztikai térkép}

\begin{longtable}{L{0.23\textwidth}L{0.30\textwidth}L{0.38\textwidth}}
\toprule
Tünet & Első hely, ahol nézd & Következő kérdés \\
\midrule
\endhead
Nem indul a service & config-check, server log & Config/secret/path/policy vagy dependency readiness? \\
\addlinespace
502/504 a proxyban & proxy log, RWLang listener/readiness & Fut-e a processz, loopback port helyes-e, readiness zöld-e? \\
\addlinespace
400/415 & access log + route body-contract & Malformed input, rossz schema vagy rossz Content-Type érkezett? \\
\addlinespace
401 & auth/session audit esemény & Nincs session, lejárt, invalidálva vagy MFA hiányzik? \\
\addlinespace
403 & route/object authorization, CSRF/CORS & Ki az actor, melyik policy tagadta meg? \\
\addlinespace
404/405/406 & route map + request method/Accept & Létezik a route, jó methodot és elfogadható reprezentációt kér a kliens? \\
\addlinespace
409 & optimistic-lock/business state & Közben megváltozott a record, vagy a \code{Changed} query nulla sort módosított? \\
\addlinespace
413/431 & request hard limitek & Body/upload vagy header lépte túl a konfigurált limitet? \\
\addlinespace
421/426 & Host/proxy/TLS policy & Helyes a public host, trusted proxy és az effektív HTTPS? \\
\addlinespace
422 & named form rerender & Melyik mező typed decode/validation hibája áll fenn? \\
\addlinespace
429 & rate-limit metrics/audit & Melyik limiter és melyik identity/IP bucket fogyott el? \\
\addlinespace
500 & request ID alapján server log & Valódi belső invariáns/internal hiba történt? \\
\addlinespace
503 & readiness + server log & DB/Redis/auth/storage vagy request resource limit tette átmenetileg kiszolgálhatatlanná? \\
\addlinespace
DB hiba & DB readiness + query környezet & Kapcsolat, timeout, schema drift, lock vagy constraint? \\
\addlinespace
Fájl/media hiba & AppFs path/jogosultság/limit & Storage root írható, confinement és byte/pixel limit rendben? \\
\addlinespace
Külső API hiba & egress policy + adapter log & DNS, CIDR, port, TLS, timeout vagy response limit zárt? \\
\addlinespace
Log eltűnik rotation után & \code{rw_log_fallback_total}, service log & SIGHUP eljutott-e, új fájl jogosultsága megfelelő-e? \\
\bottomrule
\end{longtable}

\section{Mit vigyél magaddal productionbe?}

A könyv végére a legfontosabb gyakorlati szabályok röviden:

\begin{itemize}
  \item előbb legyen typed adat és explicit policy, csak utána HTTP felület;
  \item a stabil production policy \code{server.toml}-ban legyen, ne hosszú service flaglistában;
  \item a secret érték fájlból érkezzen, ne parancssorból és ne repositoryból;
  \item schema változás explicit operátori lépés, nem szerver-startup mellékhatás;
  \item proxy header csak trusted proxy felől authority;
  \item Redis runtime/security state, nem üzleti source of truth;
  \item DB és AppFs együtt is alkothat recovery szempontból egy üzleti állapotot;
  \item a backup addig nem bizonyított, amíg izolált restore drill nem sikerült;
  \item release előtt checkelj, deploy után mérj és smoke-tesztelj;
  \item ha egy külső adat vagy erőforrás nincs explicit korlátozva, a boundary még nincs kész.
\end{itemize}
